\documentclass[10pt,letterpaper]{article}
\usepackage{cornell-notes}

\title{Clause 05.02.05: Aspects}
\author{}
\date{}
\setDocTitle{Clause 05.02.05: Aspects}
\setDocSubtitle{ISO/IEC/IEEE 42010:2022}
\setDocOwner{ISO 42010 Cornell Notes}
\setDocDate{\today}
\setCornellCollection{ISO/IEC/IEEE 42010:2022 Cornell Notes}
\setCornellUnitType{Clause}
\setCornellUnitNumber{05.02.05}
\setCornellUnitTitle{Aspects}
\hypersetup{pdftitle={Clause 05.02.05: Aspects},pdfsubject={Cornell notes},pdfauthor={}}

\begin{document}
\maketitle
\begin{CornellOverviewBox}{Overview}
{\Large\bfseries\textcolor{CNavy}{Section 5.2.5: Aspects}}\\[3pt]
{\small\textbf{Classification:} Conceptual foundation}\\
{\small\textbf{Learning objective:} Explain the role of aspects and connect it to related architecture-description concepts.}
\end{CornellOverviewBox}
\vspace{3pt}

\begin{CornellNotesTable}
\CornellNoteRow{Purpose / key concept}{Aspects capture a set of characteristics or features of the entity of interest in its environment to address concerns within an AD.}
\CornellNoteRow{Core details}{\begin{itemize}[leftmargin=*,nosep,topsep=1pt]
\item An aspect can relate to one or more concerns of stakeholders.
\item Usage of known aspects based upon prior experience within a field of application enables systematic coverage of the range of established concerns and also the identification of new concerns.
\item Aspects are based on experience in characterization of architectures and are more objective in nature as they arise from agreements among experts about practice in a domain and therefore are presumed to be best practice.
\item By examining aspects, relevant features or properties of the entity of interest can be discerned or predicted.
\end{itemize}}
\CornellNoteRow{Distinction}{An aspect is part of the entity’s character used to structure analysis; a concern is a matter of stakeholder relevance or importance.}
\CornellNoteRow{Study relationship / visual insight}{Study relationship: the entity of interest is situated in an environment; stakeholders have interests in the entity and its architecture; stakeholder perspectives give rise to concerns; aspects characterize parts of the entity’s nature and can be used to examine those concerns.}
\CornellNoteRow{Example / application}{Spatial, structural, functional, informational and programmatic aspects in an aircraft AD.}
\CornellNoteRow{Interpretive note}{\begin{itemize}[leftmargin=*,nosep,topsep=1pt]
\item A.4.2 contains more information about the utility of aspects.
\end{itemize}}
\CornellNoteRow{Related sections}{A.4.2}
\CornellNoteRow{Active recall}{\begin{itemize}[leftmargin=*,nosep,topsep=1pt]
\item What is the central role of Aspects?
\item How does Aspects connect to adjacent architecture-description concepts?
\item What closely related concepts must be kept distinct?
\end{itemize}}
\end{CornellNotesTable}


\begin{CornellSummaryBox}{Summary}
Aspects capture a set of characteristics or features of the entity of interest in its environment to address concerns within an AD. An aspect is part of the entity’s character used to structure analysis; a concern is a matter of stakeholder relevance or importance.
\end{CornellSummaryBox}

\end{document}
